\documentclass[10pt]{article}
\usepackage[letterpaper,margin=0.75in]{geometry}
\usepackage[T1]{fontenc}
\usepackage[utf8]{inputenc}
\usepackage{amsmath,amssymb}
\usepackage{array,booktabs,tabularx}
\usepackage{enumitem}
\usepackage[hidelinks]{hyperref}

\setlength{\parindent}{0pt}
\setlength{\parskip}{0.42em}
\setlist[itemize]{topsep=0.2em,itemsep=0.25em}
\newcommand{\code}[1]{\texttt{#1}}

\begin{document}
\begin{center}
{\Large\bfseries CS 224n Assignment 4}\\[0.25em]
{\large Neural Machine Translation with RNNs and Attention}
\end{center}
\vspace{0.35em}

\section{Neural Machine Translation with RNNs}

\subsection*{1(g) Encoder sentence masks}

Immediately before the softmax, every attention score corresponding to a source \code{<pad>} position is replaced by $-\infty$. Consequently, softmax assigns those positions probability zero, and their encoder states make no contribution to the attention-weighted vector $a_t$. The remaining probability mass is normalized only over real source tokens.

This is necessary because batches contain artificial padding. Without the mask, the decoder could attend to padding states, producing length-dependent noise and wasting probability mass on positions that contain no linguistic information.

\subsection*{1(i) Corpus BLEU}

Decoding the 8,064-sentence test set and scoring the translations against \code{en\_es\_data/test.en} gives:

\[
\boxed{\text{Corpus BLEU} = 21.550124}
\]

This exceeds the required score of 21. The output used for scoring is saved as \code{outputs/test\_outputs.txt}.

\subsection*{1(j) Attention mechanisms}

\begin{tabularx}{\textwidth}{>{\bfseries}p{0.17\textwidth} X X}
\toprule
Mechanism & One advantage & One disadvantage \\
\midrule
Dot product & It has no learned scoring parameters and can be implemented as a highly efficient batched matrix multiplication. & Encoder and decoder states must have compatible dimensions, and unscaled dot products can become large in high dimensions, making softmax overly sharp. \\
\addlinespace
Multiplicative & The learned matrix $W$ can align different encoder and decoder representation spaces and can also map different dimensions. & It adds parameters and matrix multiplication cost compared with plain dot-product attention. \\
\addlinespace
Additive & Its learned nonlinear scoring function is flexible and can compare states from different representation spaces; it often works well at modest hidden sizes. & It requires more parameters and per-pair operations, so it is generally slower and less parallel-friendly than dot-product attention. \\
\bottomrule
\end{tabularx}

\section{Analyzing NMT Systems}

\subsection*{2(a) Analysis of supplied translation errors}

\textbf{(i) Over-translation and repetition.}
The phrase ``another favorite of my favorites'' repeats the meaning of \emph{favorite}; the intended phrase is ``another one of my favorites.'' The attention-based recurrent decoder has no explicit coverage mechanism, so it can translate the same source content more than once. A coverage vector or coverage penalty during decoding could discourage repeated attention to already translated source positions.

\textbf{(ii) Morphology and word-order error.}
The output ``the author for children, more reading'' is ungrammatical and fails to render \emph{mas leido} as the superlative passive phrase ``most widely read.'' The model has difficulty with Spanish morphology and with a long-distance English reordering in this long sentence. Subword representations, more examples of this construction, and a stronger contextual encoder such as a Transformer would make the relevant morphological and syntactic pattern easier to learn.

\textbf{(iii) Out-of-vocabulary proper name.}
The surname ``Bolingbroke'' becomes \code{<unk>}. A fixed word-level vocabulary cannot generate a rare name that was absent or below the frequency threshold in training. Byte-pair encoding, another subword vocabulary, or character-level generation would allow the model to copy or compose unseen names.

\textbf{(iv) Lexical ambiguity and an idiom.}
The model translates \emph{manzana} as ``apple'' even though \emph{dar vuelta a la manzana} means ``go around the block.'' The word-level model selected a frequent isolated sense rather than recognizing the multiword expression. Training on more idiomatic examples, using phrase-aware contextual representations, or adding an idiom lexicon could favor the contextually correct sense.

\textbf{(v) Multiword expression and domain confusion.}
The phrase \emph{sala de profesores} should be ``teachers' lounge,'' but the model produces ``women's room.'' The decoder appears to overuse the bathroom context and fails to preserve the compound meaning of the room name. More parallel examples of institutional room names, subword/phrase-aware modeling, or terminology constraints at decoding time could correct this error.

\textbf{(vi) Numerical and world-knowledge error.}
The system changes ``hectares'' to ``acres'' but copies 100,000 instead of converting it to roughly 250,000. Standard sequence models treat numerals mostly as tokens and do not reliably perform arithmetic or unit conversion. A number-aware module or a deterministic unit-conversion postprocessor could convert the value and unit together.

\subsection*{2(b) Two additional errors from the test output}

\textbf{Example 1: dropped subject gender.}

\begin{itemize}
  \item \textbf{Source:} Lo haca en secreto,
  \item \textbf{Reference:} She did it in secret.
  \item \textbf{Model:} He did it in secret it.
  \item \textbf{Error:} The translation changes the discourse referent from female to male.
  \item \textbf{Reason:} Spanish permits omitted subjects, and the local verb form does not mark gender. The sentence-level model lacks the preceding discourse needed to recover the referent.
  \item \textbf{Possible fix:} Translate with document-level context or pass a representation of preceding sentences to the encoder so pronouns can be resolved consistently.
\end{itemize}

\textbf{Example 2: lost negative question.}

\begin{itemize}
  \item \textbf{Source:} Porque si todo est hecho de pequeas partculas y todas las pequeas partculas siguen los principios de la mecnica cuntica, entonces No debera todo seguir los principios de la mecnica cuntica?
  \item \textbf{Reference:} Because if everything is made up of little particles and all the little particles follow quantum mechanics, then shouldn't everything just follow quantum mechanics?
  \item \textbf{Model:} Because if it's all made of little particles and all the little particles follow the principles of quantum mechanics, then it should not all follow the principles of \texttt{<unk>} mechanics.
  \item \textbf{Error:} The model turns the negative rhetorical question into an awkward declarative statement and replaces ``quantum'' with an unknown token.
  \item \textbf{Reason:} The long dependency between the question marker, negation, modal verb, and subject is difficult for a recurrent decoder; attention may focus on nearby content words while dropping low-frequency functional structure.
  \item \textbf{Possible fix:} A Transformer with stronger long-range dependencies, combined with training examples that emphasize polarity and interrogative constructions, could better preserve the question's scope and subject.
\end{itemize}

These are distinct from each other and from the supplied examples: the first is a cross-sentence pronoun-resolution error, while the second is a polarity and interrogative-scope error.

\subsection*{2(c) Sentence-level BLEU}

All token counts below are case-insensitive only because every supplied sentence is already lowercased. The weights are $\lambda_1=\lambda_2=0.5$ and $\lambda_3=\lambda_4=0$.

\textbf{(i) Using both references.}

For $c_1=$ ``the love can always do,'' the matching unigrams are \{love, can, always\}, and the matching bigrams are \{love can, can always\}. Thus
\[
p_1=\frac{3}{5},\qquad p_2=\frac{2}{4}.
\]
Here $c=5$. The reference lengths are 6 and 4; both are distance 1 from $c$, so the shorter tie is selected and $r^*=4$. Since $c\ge r^*$, $BP=1$, giving
\[
\operatorname{BLEU}(c_1)=1\cdot\exp\left(\tfrac12\log\tfrac35+\tfrac12\log\tfrac24\right)
=\sqrt{0.3}=0.547723.
\]

For $c_2=$ ``love can make anything possible,'' the matching unigrams are \{love, can, anything, possible\}, and the matching bigrams are \{love can, anything possible\}. Therefore
\[
p_1=\frac{4}{5},\qquad p_2=\frac{2}{4},\qquad c=5,\qquad r^*=4,\qquad BP=1,
\]
and
\[
\operatorname{BLEU}(c_2)=\sqrt{\frac45\cdot\frac24}=\sqrt{0.4}=0.632456.
\]
BLEU ranks $c_2$ higher. I agree: $c_2$ is grammatical and preserves the meaning that love makes things possible, whereas $c_1$ ends awkwardly with ``do.''

\textbf{(ii) Using only $r_1$.}

Now $r^*=6$ for both candidates, so
\[
BP=\exp\left(1-\frac65\right)=e^{-0.2}=0.818731.
\]
The counts for $c_1$ remain $p_1=3/5$ and $p_2=2/4$, hence
\[
\operatorname{BLEU}(c_1)=0.818731\sqrt{\frac35\cdot\frac24}=0.448437.
\]
For $c_2$, only \{love, can\} match as unigrams and only \{love can\} matches as a bigram, so $p_1=2/5$ and $p_2=1/4$. Therefore
\[
\operatorname{BLEU}(c_2)=0.818731\sqrt{\frac25\cdot\frac14}=0.258905.
\]
With only $r_1$, BLEU now ranks $c_1$ higher. I disagree with that ranking: $c_2$ is still a fluent and accurate paraphrase, but the missing second reference removes the n-grams that recognize it.

\textbf{(iii) Limitation of a single reference.}
There are usually many valid translations of the same source sentence. A single reference covers only one wording, so a fluent and semantically correct paraphrase can receive a low n-gram score merely because it uses different words or syntax. This makes sentence-level scores noisy and can bias system comparisons toward the style of the lone reference.

\textbf{(iv) BLEU versus human evaluation.}

\textbf{Advantages:}
\begin{itemize}
  \item BLEU is automatic, inexpensive, and fast enough to evaluate many systems or checkpoints over a large corpus.
  \item Given fixed tokenization and references, it is reproducible and provides a standardized numerical comparison without annotator variability.
\end{itemize}

\textbf{Disadvantages:}
\begin{itemize}
  \item N-gram overlap does not directly measure meaning, factual adequacy, discourse consistency, or grammatical acceptability, so it can mis-rank valid paraphrases and serious semantic errors.
  \item BLEU is sensitive to tokenization, reference coverage, and corpus composition; a single score gives little diagnostic explanation of what a system gets wrong.
\end{itemize}

\end{document}
